\paragraph{Why budget feasibility matters.}
Budget-exhaustion events under static allocation are structural: they arise from the mismatch between local slot allocations and the global budget, not from retrieval failures or generation errors. On the 350 deep plans evaluated, 41.7\% of static executions exhaust the budget before materializing all required evidence. The flat planner eliminates these events entirely by design ($\sum_s a_s \leq B$), and the 7.7-point EM improvement confirms that the evidence acquired within the budget is more useful than the partially completed static allocation.

\paragraph{Why always-on optimization can hurt.}
The flat planner's allocation model assumes that reallocation within the global budget improves evidence quality. This assumption holds on deep plans, where multiple coupled slots compete for a budget that is locally insufficient. On shallow plans (one or two independent slots), the local allocation is already near sufficient; the planner's reallocation changes the evidence distribution in ways that can reduce answer quality, particularly on 2WikiMultiHop where the harm is concentrated ($\Delta$EM $= -0.045$ on hops-0 questions in exploratory analysis). The gate protects the shallow regime from this harm.

\paragraph{Why dependency-sensitive importance failed.}
The chain importance law ($\tau = 2(i+1) - 1$) assigns higher weights to downstream slots in the dependency chain, reflecting the intuition that later evidence acquisition steps are more important because they depend on earlier bindings. On deep plans, this produces a marginal, statistically insignificant accuracy improvement ($\Delta$EM $= +0.0086$, $p = 0.743$). The result is honest: the dependency-weighted allocator is more expensive conceptually and does not deliver measurable accuracy gains over the simpler uniform allocator. The final system uses flat allocation.

\paragraph{Why the final method is intentionally simple.}
The system components---typed evidence plans, compile-time feasibility diagnosis, budget-aware flat planner, structural-depth gate---are deliberately simple. Each component is deterministic and explainable. The feasibility diagnosis is a single arithmetic check; the gate is a threshold on a topological property; the planner is a concave allocation. There are no learned parameters, no reinforcement-learning loops, and no prompt-level heuristics in the execution path. This simplicity is a property, not a limitation: the experimental evidence supports each component, and adding complexity (e.g., dependency weighting, learned routing) has not produced measurable gains.

\paragraph{When SlotRAG is unlikely to help.}
The system's value is proportional to the prevalence and impact of budget-exhaustion events in the target workload. On workloads where plans are shallow and the budget is generous relative to slot count, budget-exhaustion events are rare and the static allocator is sufficient. On workloads where the generation model (not retrieval) is the primary bottleneck, budget-feasible planning improves completion rates but not answer accuracy. The population-level effect is bounded by the prevalence of eligible deep plans (5.4\% of the validation natural workload).

\paragraph{Heterogeneity of exploratory effects.}
The shallow-regime harm ($\Delta$EM $= -0.021$) is primarily driven by 2WikiMultiHop ($-0.045$ on hops-0 questions). On HotpotQA and MuSiQue shallow strata, the effect is near zero. This heterogeneity confirms that the gate is not universally necessary but is concentrated on the dataset where shallow budget-aware reallocation is actively harmful. The finding does not generalize across all datasets equally.

\paragraph{Honest scope of the population-level effect.}
The deep-eligible stratum constitutes 5.4\% of the natural workload. The population-level ATE of $+$0.004 EM/question reflects this low prevalence: the gate targets a small but structurally high-risk regime. The primary value is not in large population-level accuracy gains but in eliminating budget-exhaustion events and protecting shallow plans from harm. Both effects are honest and bounded.
